\documentclass{amsart}
\usepackage{amsmath,amssymb,amsthm,hyperref}
\usepackage{algorithm}
\usepackage{algpseudocode}

\newtheorem{theorem}{Theorem}
\newtheorem{lemma}{Lemma}
\newtheorem{proposition}{Proposition}
\newtheorem{corollary}{Corollary}
\theoremstyle{definition}
\newtheorem{remark}{Remark}

\DeclareMathOperator*{\argmax}{arg\,max}

\begin{document}

\title{The Legendre--Fenchel conjugate of a product of two quadratic forms}
\maketitle

\section{Setup and statement of results}

Let $A,B\in\mathbb R^{n\times n}$ be symmetric positive definite (SPD), and set
\[
q_A(x)=\tfrac12\langle Ax,x\rangle,\qquad q_B(x)=\tfrac12\langle Bx,x\rangle ,\qquad
f(x)=q_A(x)\,q_B(x).
\]
We compute the conjugate
\[
f^\ast(s)=\sup_{x\in\mathbb R^n}\bigl\{\langle s,x\rangle-f(x)\bigr\},\qquad s\in\mathbb R^n .
\]
Note that $f$ is a degree--$4$ homogeneous polynomial; it is convex iff $B$ is a positive multiple
of $A$ (see Remark~\ref{rem:nonconvex}), so in general $f^\ast$ is the conjugate of a non-convex
function.  Nevertheless $f^\ast$ has the closed evaluation procedure described below.

Throughout we use the \emph{dual norms}
\[
\|s\|_{A^{-1}}^2=\langle A^{-1}s,s\rangle,\qquad \|s\|_{B^{-1}}^2=\langle B^{-1}s,s\rangle ,
\]
which are well defined because $A,B$ are SPD (hence invertible with SPD inverses).

\begin{theorem}[Proportional case: closed single formula]\label{thm:prop}
Suppose $B=\kappa A$ for some $\kappa>0$.  Then for every $s\in\mathbb R^n$,
\begin{equation}\label{eq:prop}
f^\ast(s)=\frac34\,\Bigl(\langle A^{-1}s,s\rangle\,\langle B^{-1}s,s\rangle\Bigr)^{1/3}
        =\frac34\,\kappa^{-1/3}\,\langle A^{-1}s,s\rangle^{2/3}.
\end{equation}
\end{theorem}

\begin{theorem}[General case: evaluation via a univariate polynomial]\label{thm:gen}
Let $A,B$ be arbitrary SPD matrices and let $s\in\mathbb R^n$.  Form a Cholesky factor
$A=LL^{\!\top}$ ($L$ lower triangular, invertible), the SPD matrix
$C=L^{-1}BL^{-\top}$, an orthogonal $Q$ diagonalising $C=Q\,\mathrm{diag}(c_1,\dots,c_n)\,Q^{\!\top}$
with $c_i>0$, and the transformed right-hand side
\begin{equation}\label{eq:wdef}
w=Q^{\!\top}L^{-1}s\in\mathbb R^n,\qquad w=(w_1,\dots,w_n).
\end{equation}
Define the rational function
\begin{equation}\label{eq:gdef}
g(t)=\sum_{i=1}^n\frac{(c_i t-1)\,w_i^2}{(1+c_i t)^2},\qquad t>0 .
\end{equation}
If $w=0$ (equivalently $s=0$) then $f^\ast(s)=0$.  Otherwise $g$ has at least one zero in
$(0,\infty)$; each zero $t$ is associated with the value
\begin{equation}\label{eq:valt}
V(t)=3\,t\,\rho(t)^{2/3},\qquad
\rho(t)=\frac{1}{2t}\sum_{i=1}^n\frac{w_i^2}{(1+c_i t)^2}\;(>0),
\end{equation}
and
\begin{equation}\label{eq:fstar-gen}
f^\ast(s)=\max\bigl\{\,V(t)\ :\ t>0,\ g(t)=0\,\bigr\}.
\end{equation}
Multiplying \eqref{eq:gdef} by $\prod_{j}(1+c_j t)^2$ shows that the zeros of $g$ in $(0,\infty)$
are exactly the positive roots of the polynomial
\begin{equation}\label{eq:poly}
P(t)=\sum_{i=1}^n (c_i t-1)\,w_i^2\!\!\prod_{j\neq i}(1+c_j t)^2,\qquad \deg P\le 2n-1,
\end{equation}
so $f^\ast(s)$ is computable by finding the (finitely many) positive roots of $P$ and taking the
largest of the corresponding values $V(t)$.
\end{theorem}

Theorem~\ref{thm:prop} is the special case of Theorem~\ref{thm:gen} in which all $c_i$ are equal;
we nevertheless state and prove it separately because it yields the unconditional closed form
\eqref{eq:prop}.

\section{Reduction to a diagonal model}

\begin{lemma}[Finiteness and attainment]\label{lem:coerc}
For every $s$, the supremum defining $f^\ast(s)$ is finite and attained at some
$x^\ast\in\mathbb R^n$.  Moreover $f^\ast(s)\ge0$, with equality iff $s=0$, in which case
$x^\ast=0$.
\end{lemma}

\begin{proof}
Let $\alpha=\lambda_{\min}(A)>0$, $\beta=\lambda_{\min}(B)>0$ (positive since $A,B$ are SPD).
For all $x$, $q_A(x)\ge\frac{\alpha}{2}\|x\|^2$ and $q_B(x)\ge\frac{\beta}{2}\|x\|^2$, hence
\begin{equation}\label{eq:fcoerc}
f(x)=q_A(x)q_B(x)\ge\frac{\alpha\beta}{4}\,\|x\|^4 .
\end{equation}
Therefore $\langle s,x\rangle-f(x)\le\|s\|\,\|x\|-\frac{\alpha\beta}{4}\|x\|^4$, and the right-hand
side tends to $-\infty$ as $\|x\|\to\infty$.  In particular there is $R>0$ such that
$\langle s,x\rangle-f(x)<0=\langle s,0\rangle-f(0)$ whenever $\|x\|>R$; thus the supremum equals the
maximum of the continuous function $x\mapsto\langle s,x\rangle-f(x)$ over the compact ball
$\{\|x\|\le R\}$, which is finite and attained at some $x^\ast$.  Since $f(0)=0$ we have
$f^\ast(s)\ge\langle s,0\rangle-f(0)=0$.  If $s\ne0$, then for small $\varepsilon>0$,
$\langle s,\varepsilon s\rangle-f(\varepsilon s)=\varepsilon\|s\|^2-O(\varepsilon^4)>0$, so
$f^\ast(s)>0$; hence $x^\ast\ne0$.  If $s=0$, then by \eqref{eq:fcoerc} the objective $-f(x)\le0$
with equality only at $x=0$, so $f^\ast(0)=0$ attained uniquely at $x^\ast=0$.
\end{proof}

\begin{lemma}[Congruence reduction]\label{lem:reduce}
With $L,C,Q,w$ as in Theorem~\ref{thm:gen}, define the linear bijection
$\Phi(z)=L^{-\top}Q\,z$.  Then for $x=\Phi(z)$,
\begin{equation}\label{eq:reduce}
\langle s,x\rangle-f(x)=\langle w,z\rangle-\tfrac14\Bigl(\textstyle\sum_i z_i^2\Bigr)
\Bigl(\textstyle\sum_i c_i z_i^2\Bigr).
\end{equation}
Consequently $f^\ast(s)=\sup_{z}\bigl\{\langle w,z\rangle-\tfrac14(\sum_i z_i^2)(\sum_i c_iz_i^2)\bigr\}$.
\end{lemma}

\begin{proof}
$C=L^{-1}BL^{-\top}$ is symmetric and, for $y\ne0$, $\langle Cy,y\rangle=\langle B L^{-\top}y,L^{-\top}y\rangle>0$,
so $C$ is SPD and its eigenvalues $c_i$ are positive.  Put $y=L^{\top}x$, so $x=L^{-\top}y$ and
$q_A(x)=\tfrac12\langle LL^{\top}x,x\rangle=\tfrac12\|y\|^2$, while
$q_B(x)=\tfrac12\langle B L^{-\top}y,L^{-\top}y\rangle=\tfrac12\langle Cy,y\rangle$.  Writing
$y=Qz$ (orthogonal change of variables), $\|y\|^2=\|z\|^2=\sum_iz_i^2$ and
$\langle Cy,y\rangle=\langle Q^{\top}CQ z,z\rangle=\sum_i c_iz_i^2$.  Thus $x=L^{-\top}Qz=\Phi(z)$,
$q_A(x)=\tfrac12\sum z_i^2$, $q_B(x)=\tfrac12\sum c_iz_i^2$, so
$f(x)=\tfrac14(\sum z_i^2)(\sum c_iz_i^2)$.  Finally
$\langle s,x\rangle=\langle s,L^{-\top}Qz\rangle=\langle Q^{\top}L^{-1}s,z\rangle=\langle w,z\rangle$
by \eqref{eq:wdef}.  This is \eqref{eq:reduce}; since $\Phi$ is a bijection of $\mathbb R^n$,
taking suprema over $z$ gives the last claim.
\end{proof}

By Lemma~\ref{lem:reduce} it suffices to evaluate the conjugate of the diagonal model
\begin{equation}\label{eq:F}
F(z)=\tfrac14\Bigl(\textstyle\sum_i z_i^2\Bigr)\Bigl(\textstyle\sum_i c_iz_i^2\Bigr),
\qquad
f^\ast(s)=\sup_{z}\{\langle w,z\rangle-F(z)\}.
\end{equation}
Set $p=p(z)=\sum_iz_i^2\ge0$ and $q=q(z)=\sum_i c_iz_i^2\ge0$, so $F=\tfrac14 pq$.

\section{The stationarity system and the value $3ab$}

\begin{lemma}[Stationarity]\label{lem:stat}
Let $z^\ast$ be a global maximiser of \eqref{eq:F}.  Put
$a=q_{A}(\Phi(z^\ast))=\tfrac12 p(z^\ast)$ and $b=q_{B}(\Phi(z^\ast))=\tfrac12 q(z^\ast)$.
Then
\begin{equation}\label{eq:zform}
w_i=z_i^\ast\,(b+a\,c_i)\qquad(i=1,\dots,n),
\end{equation}
and
\begin{equation}\label{eq:value3ab}
f^\ast(s)=\langle w,z^\ast\rangle-F(z^\ast)=3ab .
\end{equation}
If $s\neq0$ then $a>0$, $b>0$, $b+ac_i>0$ for all $i$, and $z_i^\ast=w_i/(b+ac_i)$.
\end{lemma}

\begin{proof}
$F$ is a polynomial, hence smooth, and $z\mapsto\langle w,z\rangle-F(z)$ attains its global maximum
at $z^\ast$ (Lemma~\ref{lem:coerc}, transported by the bijection $\Phi$).  Hence
$\nabla_z\bigl(\langle w,z\rangle-F(z)\bigr)=0$ at $z^\ast$:
\[
w_i=\partial_{z_i}F(z^\ast)
   =\tfrac14\Bigl(2z_i^\ast\,q(z^\ast)+p(z^\ast)\,2c_iz_i^\ast\Bigr)
   =\tfrac12 z_i^\ast\bigl(q(z^\ast)+c_i\,p(z^\ast)\bigr)
   = z_i^\ast\,(b+a\,c_i),
\]
using $a=\tfrac12 p$, $b=\tfrac12 q$.  This is \eqref{eq:zform}.

Since $A,B$ are SPD we have $a=\tfrac12 p(z^\ast)\ge0$ and $b=\tfrac12 q(z^\ast)\ge0$, with
$c_i>0$.  Assume now $s\ne0$.  By Lemma~\ref{lem:coerc} (transported by $\Phi$) $z^\ast\ne0$, so
$p(z^\ast)>0$ and, since all $c_i>0$, also $q(z^\ast)>0$; thus $a>0$ and $b>0$.  Consequently
$b+ac_i\ge b>0$ for every $i$, and \eqref{eq:zform} gives $z_i^\ast=w_i/(b+ac_i)$.

Finally, from \eqref{eq:zform} (valid for all $s$),
\[
\langle w,z^\ast\rangle=\sum_i w_iz_i^\ast=\sum_i z_i^{\ast2}(b+ac_i)
=b\sum_i z_i^{\ast2}+a\sum_ic_iz_i^{\ast2}
=b\,p(z^\ast)+a\,q(z^\ast)=2ab+2ab=4ab,
\]
where we used $p(z^\ast)=2a$ and $q(z^\ast)=2b$.  Also $F(z^\ast)=\tfrac14 p q=\tfrac14(2a)(2b)=ab$.
Hence the optimal value equals $\langle w,z^\ast\rangle-F(z^\ast)=4ab-ab=3ab$, proving
\eqref{eq:value3ab}.  (When $s=0$, $z^\ast=0$ gives $a=b=0$ and $3ab=0=f^\ast(0)$, consistent with
Lemma~\ref{lem:coerc}.)
\end{proof}

\section{Proof of Theorem~\ref{thm:prop} (proportional case)}

\begin{proof}
If $s=0$ both sides of \eqref{eq:prop} vanish (Lemma~\ref{lem:coerc}), so assume $s\neq0$.
Because $B=\kappa A$ with $\kappa>0$, $f(x)=q_A(x)\,\kappa q_A(x)=\tfrac\kappa4\langle Ax,x\rangle^2$,
and
\[
f^\ast(s)=\sup_{x}\Bigl\{\langle s,x\rangle-\tfrac\kappa4\langle Ax,x\rangle^2\Bigr\}.
\]
Fix the value $u:=\langle Ax,x\rangle\ge0$ and maximise the linear term over $\{x:\langle Ax,x\rangle=u\}$.
By Cauchy--Schwarz in the inner product $\langle A\cdot,\cdot\rangle$,
$\langle s,x\rangle=\langle A^{-1}s,Ax\rangle\le\sqrt{\langle A^{-1}s,s\rangle}\,\sqrt{\langle Ax,x\rangle}
=\|s\|_{A^{-1}}\sqrt{u}$,
with equality for a suitable $x$ (namely $x$ a nonnegative multiple of $A^{-1}s$, which realises
every $u\ge0$).  Hence
\[
f^\ast(s)=\sup_{u\ge0}\Bigl\{\|s\|_{A^{-1}}\sqrt{u}-\tfrac\kappa4 u^2\Bigr\}.
\]
Let $N=\|s\|_{A^{-1}}>0$ and $h(u)=N\sqrt u-\tfrac\kappa4u^2$.  On $(0,\infty)$,
$h'(u)=\tfrac{N}{2\sqrt u}-\tfrac\kappa2 u$, which is strictly decreasing (since
$\frac{N}{2\sqrt u}$ is decreasing and $-\frac\kappa2u$ is decreasing), with $h'(0^+)=+\infty$ and
$h'(+\infty)=-\infty$; thus $h$ has a unique critical point $u^\ast>0$, which is its global
maximiser, solving $\tfrac{N}{2\sqrt{u^\ast}}=\tfrac\kappa2 u^\ast$, i.e.\
$N=\kappa\,u^{\ast 3/2}$, so $u^\ast=(N/\kappa)^{2/3}$.  Therefore
\[
f^\ast(s)=h(u^\ast)=N\,(N/\kappa)^{1/3}-\tfrac\kappa4(N/\kappa)^{4/3}
=N^{4/3}\kappa^{-1/3}\Bigl(1-\tfrac14\Bigr)=\tfrac34\,\kappa^{-1/3}N^{4/3}.
\]
Since $N^{4/3}=\langle A^{-1}s,s\rangle^{2/3}$, this is the second equality in \eqref{eq:prop}.
For the first, $B=\kappa A$ gives $B^{-1}=\kappa^{-1}A^{-1}$, hence
$\langle B^{-1}s,s\rangle=\kappa^{-1}\langle A^{-1}s,s\rangle$, so
$\bigl(\langle A^{-1}s,s\rangle\langle B^{-1}s,s\rangle\bigr)^{1/3}
=\kappa^{-1/3}\langle A^{-1}s,s\rangle^{2/3}$, which agrees.
\end{proof}

\begin{remark}\label{rem:nonconvex}
$f$ is convex iff $B=\kappa A$ for some $\kappa>0$, which (via Lemma~\ref{lem:reduce}) is equivalent
to all $c_i$ being equal.  If $B=\kappa A$ then $f=\tfrac\kappa4\langle A\cdot,\cdot\rangle^2$ is a
nondecreasing convex function ($u\mapsto\tfrac\kappa4u^2$ on $u\ge0$) of the convex nonnegative map
$x\mapsto\langle Ax,x\rangle$, hence convex.  Conversely, suppose not all $c_i$ are equal, say
$c_1\ne c_2$; by scaling we may take $c_1=1$.  The $2\times2$ leading block of the Hessian of $F$ at
the point $z=2e_1+e_2$ (i.e.\ $z_1=2,z_2=1$, other coordinates $0$) is, by direct differentiation of
$F=\tfrac14(\sum z_i^2)(\sum c_iz_i^2)$,
\[
\begin{pmatrix}
\partial^2_{z_1z_1}F & \partial^2_{z_1z_2}F\\[2pt]
\partial^2_{z_2z_1}F & \partial^2_{z_2z_2}F
\end{pmatrix}
=\begin{pmatrix} \tfrac{15}{2}+\tfrac{c_2}{2} & 2(1+c_2)\\[2pt] 2(1+c_2) & \tfrac{5}{2}+\tfrac{5c_2}{2}\end{pmatrix},
\]
whose determinant equals $-\tfrac34\,(c_2-1)^2<0$ (using $c_2\ne1$).  A symmetric matrix with a
$2\times2$ principal submatrix of negative determinant is indefinite, so the Hessian of $F$ is not
positive semidefinite at this point; hence $F$, and therefore $f=F\circ\Phi^{-1}$ up to the linear
change of variables, is non-convex.  Thus a single unconditional closed formula \eqref{eq:prop} is
available exactly in the proportional case.
\end{remark}

\section{Proof of Theorem~\ref{thm:gen} (general case)}

\begin{proof}
We use the diagonal model \eqref{eq:F}.  If $s=0$ then $w=0$ and $f^\ast(s)=0$ by
Lemma~\ref{lem:coerc}; assume henceforth $s\ne0$, so $w=Q^{\top}L^{-1}s\ne0$ (as $Q^{\top}L^{-1}$ is
invertible).

\subsection*{Step 1: positive scaling of the stationarity system}
By Lemma~\ref{lem:stat} there is a global maximiser $z^\ast$ with $a>0,b>0$, $b+ac_i>0$, and
$z_i^\ast=w_i/(b+ac_i)$, and $f^\ast(s)=3ab$.  Substituting $z_i^\ast=w_i/(b+ac_i)$ into
$a=\tfrac12 p(z^\ast)=\tfrac12\sum_i z_i^{\ast2}$ and $b=\tfrac12 q(z^\ast)=\tfrac12\sum_i c_iz_i^{\ast2}$
gives the closed two--equation system in $(a,b)$:
\begin{equation}\label{eq:ab-sys}
a=\frac12\sum_{i=1}^n\frac{w_i^2}{(b+ac_i)^2},\qquad
b=\frac12\sum_{i=1}^n\frac{c_i\,w_i^2}{(b+ac_i)^2}.
\end{equation}
Introduce the ratio $t=a/b>0$.  Writing $a=tb$, we have $b+ac_i=b(1+c_i t)$, so \eqref{eq:ab-sys}
becomes
\begin{equation}\label{eq:scaled}
tb=\frac{1}{2b^2}\sum_i\frac{w_i^2}{(1+c_it)^2},\qquad
b=\frac{1}{2b^2}\sum_i\frac{c_i w_i^2}{(1+c_it)^2}.
\end{equation}
Multiplying the first equation by $b^2/t$ and the second by $b^2$ yields, with $\rho(t)$ as in
\eqref{eq:valt},
\begin{equation}\label{eq:b3}
b^3=\frac{1}{2t}\sum_i\frac{w_i^2}{(1+c_it)^2}=\rho(t),\qquad
b^3=\frac12\sum_i\frac{c_iw_i^2}{(1+c_it)^2}.
\end{equation}

\subsection*{Step 2: the univariate equation $g(t)=0$ and its associated value $V(t)$}
\emph{(Necessity.)} Equating the two expressions for $b^3$ in \eqref{eq:b3} and multiplying by
$2t>0$,
\[
\sum_i\frac{w_i^2}{(1+c_it)^2}=t\sum_i\frac{c_iw_i^2}{(1+c_it)^2}
\iff
\sum_i\frac{(c_it-1)\,w_i^2}{(1+c_it)^2}=0,
\]
i.e.\ $g(t)=0$ with $g$ as in \eqref{eq:gdef}.  Thus the ratio $t^\ast=a/b$ of the global maximiser
is a zero of $g$ in $(0,\infty)$; in particular such a zero exists.  Moreover, by
\eqref{eq:b3}, $b^2=\rho(t^\ast)^{2/3}$, and by Lemma~\ref{lem:stat},
\[
f^\ast(s)=3ab=3(t^\ast b)b=3t^\ast b^2=3t^\ast\rho(t^\ast)^{2/3}=V(t^\ast).
\]

\emph{(Each positive zero gives an attained value.)} Conversely, let $t>0$ satisfy $g(t)=0$.  Since
$w\neq0$ and every denominator $(1+c_it)^2>0$, the sum $\sum_i w_i^2/(1+c_it)^2$ is strictly
positive (terms with $w_i=0$ vanish, but not all $w_i$ are zero), so $\rho(t)>0$.  Set
$b:=\rho(t)^{1/3}>0$ and $a:=t\,b>0$.  Equating $g(t)=0$ shows the two right-hand sides in
\eqref{eq:b3} are equal, both equal to $\rho(t)=b^3$, so $(a,b)$ solves \eqref{eq:scaled}, hence
\eqref{eq:ab-sys}.  Define $z_i:=w_i/(b+ac_i)$ (legitimate since $b+ac_i\ge b>0$).  Then
$\tfrac12\sum_i z_i^2=a$ and $\tfrac12\sum_i c_iz_i^2=b$ by \eqref{eq:ab-sys}, i.e.\
$p(z)=2a$, $q(z)=2b$, and $w_i=z_i(b+ac_i)=\tfrac12 z_i(q(z)+c_ip(z))=\partial_{z_i}F(z)$.  Hence $z$
is a stationary point of the smooth objective $z\mapsto\langle w,z\rangle-F(z)$.  As in
Lemma~\ref{lem:stat}, its objective value is
\[
\langle w,z\rangle-F(z)=4ab-ab=3ab=3t\,b^2=3t\,\rho(t)^{2/3}=V(t),
\]
which is a value attained by the objective.

\subsection*{Step 3: the conjugate equals the largest associated value}
By Step~2 there is at least one zero $t>0$ of $g$.  For every such zero, $V(t)$ is a value attained
by the objective $z\mapsto\langle w,z\rangle-F(z)$, hence $V(t)\le f^\ast(s)$.  On the other hand,
the global maximum $f^\ast(s)$ equals $V(t^\ast)$ for the maximiser's ratio $t^\ast$, which is one of
the zeros.  Therefore
\[
f^\ast(s)=\max\{V(t):t>0,\ g(t)=0\},
\]
which is \eqref{eq:fstar-gen}.  The maximum is over a finite set: multiplying $g(t)=0$ by
$\prod_j(1+c_jt)^2>0$ shows the positive zeros of $g$ coincide with the positive roots of the
polynomial $P$ of \eqref{eq:poly}.  Its degree is at most $2n-1$ (the term $i$ contributes degree
$1+2(n-1)=2n-1$), so $P$ has finitely many roots.  This proves Theorem~\ref{thm:gen}.
\end{proof}

\section{Algorithm}

\begin{algorithm}[h]
\caption{Evaluation of $f^\ast(s)$ for $f=q_Aq_B$, $A,B$ SPD}
\begin{algorithmic}[1]
\Require SPD matrices $A,B\in\mathbb R^{n\times n}$; vector $s\in\mathbb R^n$.
\Ensure The value $f^\ast(s)=\sup_x\{\langle s,x\rangle-q_A(x)q_B(x)\}$.
\If{$s=0$} \Return $0$ \EndIf
\State Compute a Cholesky factor $L$ with $A=LL^{\top}$.
\State Form $C\gets L^{-1}BL^{-\top}$ and its spectral decomposition
       $C=Q\,\mathrm{diag}(c_1,\dots,c_n)\,Q^{\top}$ with $c_i>0$.
\State Set $w\gets Q^{\top}L^{-1}s$.
\If{all $c_i$ are equal to a common value $\kappa$ \textbf{(proportional case)}}
   \State \Return $\tfrac34\,(\langle A^{-1}s,s\rangle\,\langle B^{-1}s,s\rangle)^{1/3}$
          \Comment{Theorem~\ref{thm:prop}}
\EndIf
\State Form $P(t)=\sum_{i=1}^n (c_i t-1)w_i^2\prod_{j\neq i}(1+c_j t)^2$ and compute its real roots.
\State $\mathcal R\gets\{\text{roots }t\text{ of }P\text{ with } t>0\}$.
\State $\mathrm{best}\gets -\infty$.
\For{$t\in\mathcal R$}
   \State $\rho\gets \dfrac{1}{2t}\sum_{i=1}^n \dfrac{w_i^2}{(1+c_i t)^2}$;\quad
          $V\gets 3\,t\,\rho^{2/3}$.
   \If{$V>\mathrm{best}$} $\mathrm{best}\gets V$ \EndIf
\EndFor
\State \Return $\mathrm{best}$.
\end{algorithmic}
\end{algorithm}

\paragraph{Correctness.}
Lines~1 and 6 are justified by Lemma~\ref{lem:coerc} and Theorem~\ref{thm:prop} respectively.  In
the general branch, Lemma~\ref{lem:reduce} reduces the problem to the diagonal model \eqref{eq:F},
and Theorem~\ref{thm:gen} (formulas \eqref{eq:fstar-gen}--\eqref{eq:poly}) shows that
$f^\ast(s)=\max_{t\in\mathcal R}V(t)$, which is exactly what the loop computes.  The set
$\mathcal R$ is nonempty because $g$ (hence $P$) has a positive root (Step~2 of the proof of
Theorem~\ref{thm:gen}).

\paragraph{Termination and complexity.}
The Cholesky factorisation, the $n\times n$ symmetric eigendecomposition, and the matrix products
cost $O(n^3)$.  Building $P$ (a polynomial of degree $\le 2n-1$) and finding its real roots cost
$O(n^3)$ by the companion-matrix eigenvalue method.  The final loop runs over at most $2n-1$ roots,
each iteration costing $O(n)$.  Hence the algorithm terminates in $O(n^3)$ arithmetic operations.

\begin{remark}
$f^\ast$ is positively homogeneous of degree $4/3$: from $f(\lambda x)=\lambda^4 f(x)$ one gets, for
$\mu>0$, $f^\ast(\mu s)=\sup_x\{\mu\langle s,x\rangle-f(x)\}=\sup_y\{\mu^{4/3}(\langle s,y\rangle-f(y))\}
=\mu^{4/3}f^\ast(s)$ via $x=\mu^{1/3}y$.  This is consistent with \eqref{eq:prop} and with the fact
that, in \eqref{eq:valt}, $w\mapsto\mu w$ leaves the zeros of $g$ unchanged while multiplying
$\rho(t)$ by $\mu^2$, hence $V(t)$ by $\mu^{4/3}$.
\end{remark}

\end{document}
